\documentclass[10pt,a4paper]{article}
\usepackage[utf8]{inputenc}
\usepackage[vietnamese]{babel}
\usepackage{amsmath,amssymb,amsfonts}
\usepackage{geometry}
\usepackage{xcolor}
\usepackage{tcolorbox}
\usepackage{fancyhdr}
\usepackage{mathptmx}

\geometry{
    a4paper,
    left=62mm,
    right=20mm,
    top=22mm,
    bottom=22mm,
    headheight=14pt,
    footskip=12mm
}

% Định nghĩa màu sắc theo chuẩn Stewart Calculus
\definecolor{stewartcyan}{RGB}{0, 136, 170}
\definecolor{stewartred}{RGB}{204, 30, 30}

\newcommand{\stewarttombstone}{\textcolor{stewartcyan}{\rule{1.1ex}{1.1ex}}}

% Khung Quy tắc 4
\newtcolorbox{rulebox}{
    colback=white,
    colframe=stewartred,
    boxrule=0.8pt,
    arc=0mm,
    left=8pt,
    right=8pt,
    top=7pt,
    bottom=7pt
}

% Thiết lập tiêu đề đầu trang dóng rộng sang lề trái
\pagestyle{fancy}
\fancyhf{}
\fancyhead[L]{\hspace{-42mm}\textsf{\textbf{202}}\qquad \textsf{\textbf{CHƯƠNG 3}}\quad \textsf{Các quy tắc đạo hàm}}
\renewcommand{\headrulewidth}{0pt}

\setlength{\parindent}{0pt}
\setlength{\parskip}{5pt}

\begin{document}

Trong Ví dụ 2(a), chúng ta đã kết hợp Quy tắc chuỗi với quy tắc tính đạo hàm của hàm sin. Tổng quát, nếu $y = \sin u$, trong đó $u$ là một hàm khả vi theo $x$, thì theo Quy tắc chuỗi,
\[
\frac{dy}{dx} = \frac{dy}{du}\,\frac{du}{dx} = \cos u\,\frac{du}{dx}
\]

\noindent\rlap{Do đó,}\hfill $\displaystyle \frac{d}{dx}(\sin u) = \cos u\,\frac{du}{dx}$\hfill\null

\vspace{3pt}
Tương tự như vậy, tất cả các công thức tính đạo hàm của các hàm lượng giác đều có thể kết hợp với Quy tắc chuỗi.

Chúng ta hãy làm rõ trường hợp đặc biệt của Quy tắc chuỗi khi hàm bên ngoài $f$ là một hàm lũy thừa. Nếu $y = [g(x)]^n$, ta có thể viết $y = f(u) = u^n$ với $u = g(x)$. Bằng cách áp dụng Quy tắc chuỗi và sau đó là Quy tắc lũy thừa, ta được
\[
\frac{dy}{dx} = \frac{dy}{du}\,\frac{du}{dx} = n u^{n-1}\,\frac{du}{dx} = n [g(x)]^{n-1} g'(x)
\]

\begin{rulebox}
    \noindent
    {\setlength{\fboxsep}{2pt}\setlength{\fboxrule}{0.8pt}\textcolor{stewartred}{\fbox{\bfseries\textsf{\color{stewartred}4}}}}\;\;
    \textbf{\textsf{\color{stewartred}Quy tắc lũy thừa kết hợp với Quy tắc chuỗi}}\quad
    Nếu $n$ là một số thực bất kỳ và $u = g(x)$ là hàm khả vi, thì
    \[
    \frac{d}{dx}(u^n) = n u^{n-1}\,\frac{du}{dx}
    \]
    \noindent\rlap{Cách khác,}\hfill $\displaystyle \frac{d}{dx}[g(x)]^n = n [g(x)]^{n-1} g'(x)$\hfill\null
\end{rulebox}

\vspace{3pt}
Lưu ý rằng đạo hàm mà chúng ta đã tìm thấy trong Ví dụ 1 có thể được tính bằng cách lấy $n = \frac{1}{2}$ trong Quy tắc 4.

\vspace{8pt}
\noindent
\textbf{\textsf{\color{stewartcyan}VÍ DỤ 3}}\quad Tính đạo hàm của $y = (x^3 - 1)^{100}$.

\vspace{3pt}
\noindent
\textbf{\textsf{\color{stewartcyan}LỜI GIẢI}}\quad Đặt $u = g(x) = x^3 - 1$ và $n = 100$ vào (4), ta có
\begin{align*}
\frac{dy}{dx} &= \frac{d}{dx}(x^3 - 1)^{100} = 100(x^3 - 1)^{99}\,\frac{d}{dx}(x^3 - 1) \\[4pt]
&= 100(x^3 - 1)^{99} \cdot 3x^2 = 300x^2(x^3 - 1)^{99} \tag*{\stewarttombstone}
\end{align*}

\vspace{8pt}
\noindent
\textbf{\textsf{\color{stewartcyan}VÍ DỤ 4}}\quad Tìm $f'(x)$ nếu $f(x) = \dfrac{1}{\sqrt[3]{x^2 + x + 1}}$.

\vspace{3pt}
\noindent
\textbf{\textsf{\color{stewartcyan}LỜI GIẢI}}\quad Trước hết viết lại $f$:\qquad $f(x) = (x^2 + x + 1)^{-1/3}$

\begin{align*}
\text{Do đó,}\qquad f'(x) &= -\frac{1}{3}(x^2 + x + 1)^{-4/3}\,\frac{d}{dx}(x^2 + x + 1) \\[4pt]
&= -\frac{1}{3}(x^2 + x + 1)^{-4/3}(2x + 1) \tag*{\stewarttombstone}
\end{align*}

\end{document}